\chapter{Nền Tảng Lý Thuyết Số \& Cơ Sở Toán Học Mật Mã Khóa Công Khai}

\banddo{Nền tảng đại số trừu tượng $\to$ Bài toán khó (Factorization, DLP, ECDLP) $\to$ Cơ chế Trapdoor $\to$ Bài toán mạo danh MitM $\to$ Định danh Identity Binding}

\section{Cấu Trúc Đại Số Trừu Tượng \& Các Phép Toán Số Học}

Mật mã khóa công khai không thể tồn tại nếu thiếu nền tảng vững chắc của đại số trừu tượng. Thay vì hoạt động trên tập số thực liên tục $\mathbb{R}$, mật mã học hiện đại vận hành trên các cấu trúc rời rạc hữu hạn nhằm bảo đảm tính chính xác tuyệt đối và ngăn chặn triệt để hiện tượng sai số làm tròn số học.

\subsection{Nhóm (Groups), Vành (Rings) và Trường Hữu Hạn (Finite Fields)}

\begin{definition}[Nhóm đại số]
Một tập hợp $G$ cùng với phép toán hai ngôi $\star$ được gọi là một \emph{nhóm} $(G, \star)$ nếu thỏa mãn 4 tiên đề:
\begin{enumerate}
  \item \textbf{Tính đóng (Closure):} $\forall a, b \in G \implies a \star b \in G$.
  \item \textbf{Tính kết hợp (Associativity):} $\forall a, b, c \in G \implies (a \star b) \star c = a \star (b \star c)$.
  \item \textbf{Phần tử trung hòa (Identity):} $\exists e \in G, \forall a \in G \implies a \star e = e \star a = a$.
  \item \textbf{Phần tử nghịch đảo (Inverse):} $\forall a \in G, \exists a^{-1} \in G \implies a \star a^{-1} = a^{-1} \star a = e$.
\end{enumerate}
Nếu phép toán có tính giao hoán $a \star b = b \star a$, nhóm được gọi là \emph{nhóm Abel}.
\Enddinhnghia
\end{definition}

Vành số dư $\mathbb{Z}_n = \{0, 1, \dots, n-1\}$ là vành cơ bản của hệ mật RSA. Trường hữu hạn Galois $\mathbb{F}_q$ chỉ tồn tại khi số phần tử $q = p^m$ (với $p$ là số nguyên tố và $m \ge 1$):
\begin{itemize}
  \item Khi $m = 1$, ta có trường số nguyên tố $\mathbb{F}_p = \mathbb{Z}_p$, nơi các phép toán modulo $p$ diễn ra.
  \item Khi $m > 1$, ta có trường mở rộng $\mathbb{F}_{2^m}$, đặc biệt quan trọng trong các bộ đồng xử lý mật mã phần cứng.
\end{itemize}

\subsection{Các Định Lý Số Học Cốt Lõi}

\begin{theorem}[Bổ đề Bézout \& Nghịch đảo Modulo]
Với hai số nguyên $a, b$, luôn tồn tại hai số nguyên $x, y$ sao cho:
\begin{equation}
a \cdot x + b \cdot y = \gcd(a, b)
\end{equation}
Khi $\gcd(a, n) = 1$, ta có $a \cdot x \equiv 1 \pmod n$, tức $x = a^{-1} \pmod n$. Thuật toán Euclid mở rộng tính toán phần tử nghịch đảo này trong thời gian đa thức cực nhanh $\mathcal{O}(\log n)$.
\Endtheorem
\end{theorem}

\begin{theorem}[Định lý Euler \& Định lý Fermat nhỏ]
Với số nguyên dương $n$, hàm phi Euler $\phi(n)$ đếm số lượng các số nguyên dương nhỏ hơn $n$ và nguyên tố cùng nhau với $n$. Nếu $n = p \cdot q$ ($p, q$ nguyên tố), thì $\phi(n) = (p - 1)(q - 1)$. Khi đó, với mọi $a$ nguyên tố cùng nhau với $n$:
\begin{equation}
a^{\phi(n)} \equiv 1 \pmod n
\end{equation}
Trường hợp đặc biệt khi $n = p$ là số nguyên tố, ta thu được Định lý Fermat nhỏ: $a^{p-1} \equiv 1 \pmod p$.
\Endtheorem
\end{theorem}

\begin{lichsu}
Năm 1976, hai nhà nghiên cứu tại Đại học Stanford là Whitfield Diffie và Martin Hellman đã công bố bài báo lịch sử \emph{''New Directions in Cryptography''} (kế thừa các ý tưởng Merkle Puzzles của Ralph Merkle). Công trình này đã phá vỡ thế độc tôn hàng nghìn năm của mật mã đối xứng, đặt nền móng cho cuộc cách mạng mật mã khóa công khai. Chỉ một năm sau đó (1977), Ron Rivest, Adi Shamir và Leonard Adleman tại MIT đã hiện thực hóa khái niệm này bằng thuật toán RSA lừng danh.
\end{lichsu}

\section{Hàm Một Chiều Có Cửa Sập \& Bài Toán Khó}

Mọi hệ thống mật mã khóa công khai đều được xây dựng trên nền tảng của \textbf{Hàm một chiều có cửa sập (Trapdoor One-Way Function -- TOWF)}:
\begin{equation}
f: X \to Y
\end{equation}
Hàm $f$ phải thỏa mãn ba đặc trưng toán học nghiêm ngặt:
\begin{enumerate}
  \item \textbf{Dễ tính toán theo chiều xuôi:} Cho trước $x \in X$, việc tính $y = f(x)$ là một bài toán thời gian đa thức $\mathcal{O}(\text{poly}(|x|))$.
  \item \textbf{Bất khả thi theo chiều ngược (Hard to Invert):} Cho trước $y \in Y$, việc tìm $x$ sao cho $f(x) = y$ là bài toán bất khả thi về mặt độ phức tạp tính toán đối với bất kỳ thuật toán cổ điển nào.
  \item \textbf{Cửa sập bí mật (Trapdoor Key):} Tồn tại một thông tin bổ trợ đặc biệt $K_{\text{priv}}$. Khi biết $K_{\text{priv}}$, việc tính ngược $x = f^{-1}(y, K_{\text{priv}})$ trở thành một bài toán đa thức tường minh và hiệu quả.
\end{enumerate}

\begin{batngo}
Toán học mật mã thuần túy rất mạnh mẽ, nhưng \textbf{nó không thể tự giải quyết được bài toán lòng tin}. Kẻ tấn công đứng giữa kênh truyền (Mallory) hoàn toàn có thể thay thế khóa công khai của Alice bằng khóa của chính hắn mà Bob không có cách nào phát hiện nếu chỉ dựa vào số học. Để giải quyết nguy cơ này, năm 1978, chàng sinh viên cử nhân Loren Kohnfelder tại MIT đã đề xuất khái niệm \emph{''Digital Certificate''} (chứng chỉ số) --- gắn chặt khóa công khai với định danh pháp lý thông qua chữ ký của bên thứ ba đáng tin cậy.
\end{batngo}

\begin{chandung}{Loren Kohnfelder}
Loren Kohnfelder hoàn thành luận văn cử nhân tại MIT vào tháng 5/1978 dưới sự hướng dẫn của Giáo sư Leonard Adleman (chữ 'A' trong RSA). Bản luận văn mang tên \emph{''Towards a Practical Public-Key Cryptosystem''} lần đầu tiên trong lịch sử nhân loại định nghĩa cấu trúc của một chứng chỉ số được ký bởi một nhà cầm quyền trung ương (Central Authority / CA). Cấu trúc mà Kohnfelder phác thảo năm 1978 chính là tiền thân trực tiếp của chuẩn quốc tế ITU-T X.509 thống trị toàn bộ Internet ngày nay.
\end{chandung}

\section{Hệ Mật RSA \& Tối Ưu Hóa RSA-CRT}

Quy trình sinh khóa RSA:
\begin{enumerate}
  \item Chọn ngẫu nhiên hai số nguyên tố rất lớn $p, q$ ($p \ne q$, thường có độ dài 1024-bit hoặc 2048-bit).
  \item Tính $n = p \cdot q$ và $\phi(n) = (p - 1)(q - 1)$.
  \item Chọn số mũ công khai $e$ thỏa mãn $\gcd(e, \phi(n)) = 1$ (tiêu chuẩn quốc tế ấn định $e = 65537 = 2^{16} + 1$).
  \item Tính số mũ bí mật $d \equiv e^{-1} \pmod{\phi(n)}$ bằng thuật toán Euclid mở rộng.
\end{enumerate}

Kỹ thuật tăng tốc RSA-CRT: Thay vì tính $S = M^d \pmod n$ trực tiếp, ta tính trên hai nửa modulo $p$ và $q$:
\begin{align}
S_p &= M^{d \pmod{p-1}} \pmod p\\
S_q &= M^{d \pmod{q-1}} \pmod q\\
h &= q^{-1} \cdot (S_p - S_q) \pmod p\\
S &= S_q + h \cdot q
\end{align}
Kỹ thuật này giúp giảm độ phức tạp tính toán và tăng tốc độ ký số thực tế lên \textbf{4 lần}.

\section{Mật Mã Đường Cong Elliptic (ECC)}

Đường cong Elliptic trên trường hữu hạn $\mathbb{F}_p$ ($p > 3$) dạng rút gọn Weierstrass:
\begin{equation}
E: y^2 = x^3 + a \cdot x + b \pmod p \quad \text{với } 4a^3 + 27b^2 \not\equiv 0 \pmod p
\end{equation}

Tập hợp các điểm $(x, y) \in \mathbb{F}_p \times \mathbb{F}_p$ thỏa mãn phương trình cùng với điểm vô cực $\mathcal{O}$ tạo thành một nhóm Abel hữu hạn dưới phép cộng điểm hình học (Chord-and-Tangent).

\begin{table}[h]
\centering\small
\begin{tabularx}{\linewidth}{@{}l c c c c@{}}
\toprule
\textbf{Mức An Toàn} & \textbf{Khóa Đối Xứng} & \textbf{Khóa RSA} & \textbf{Khóa ECC} & \textbf{Tỉ Lệ Tiết Kiệm} \\
\midrule
112-bit & 3TDEA & 2048 bits & 224 bits & 9.1 : 1 \\
\textbf{128-bit (Chuẩn hiện đại)} & \textbf{AES-128} & \textbf{3072 bits} & \textbf{256 bits} & \textbf{12.0 : 1} \\
\textbf{192-bit (An ninh cao)} & \textbf{AES-192} & \textbf{7680 bits} & \textbf{384 bits} & \textbf{20.0 : 1} \\
\textbf{256-bit (Tuyệt mật)} & \textbf{AES-256} & \textbf{15360 bits} & \textbf{521 bits} & \textbf{29.5 : 1} \\
\bottomrule
\end{tabularx}
\caption{Đối chiếu mức độ an toàn theo tiêu chuẩn NIST SP 800-57 Part 1 Rev 5.}
\end{table}

\begin{chayngay}
Sinh cặp khóa ECC Ed25519 hiện đại siêu nhanh bằng OpenSSL 3.5:
\begin{lstlisting}[language=bash]
# Sinh khoa bi mat Ed25519
openssl genpkey -algorithm ED25519 -out ed25519.key
# Trich xuat khoa cong khai tuong ung
openssl pkey -in ed25519.key -pubout -out ed25519.pub
# In chi tiet thong so khoa
openssl pkey -in ed25519.key -text -noout
\end{lstlisting}
\end{chayngay}

\begin{thuthach}
\textbf{Cạm bẫy tấn công Bleichenbacher's Oracle trên RSA PKCS\#1 v1.5:}
Trong giao thức TLS cũ sử dụng RSA trao đổi khóa, nếu máy chủ trả về các thông điệp lỗi khác nhau khi giải mã đệm PKCS\#1 v1.5 sai, kẻ tấn công có thể gửi hàng nghìn truy vấn thích nghi để giải mã toàn bộ thông điệp mà không cần khóa bí mật. Tiêu chuẩn RFC 5246 và TLS 1.3 đã loại bỏ hoàn toàn cơ chế RSA key exchange và chuyển sang Ephemeral Diffie-Hellman (ECDHE).
\end{thuthach}

\begin{baitoanmo}
Năm 1994, Peter Shor công bố thuật toán lượng tử chứng minh rằng cả bài toán phân tích thừa số nguyên tố (nền tảng RSA) và logarit rời rạc (nền tảng DH/ECC) đều bị giải quyết hoàn toàn trong thời gian đa thức $\mathcal{O}((\log N)^3)$. Khi máy tính lượng tử quy mô lớn (CRQC) ra đời, toàn bộ hệ mật RSA và ECC sẽ sụp đổ. Đây là động lực của cuộc đại di chuyển sang Mật mã Hậu lượng tử (PQC).
\end{baitoanmo}

\begin{tukiemtra}
\begin{enumerate}
  \item Tại sao thuật toán Euclid mở rộng lại đóng vai trò sống còn trong việc sinh khóa bí mật của hệ mật RSA?
  \item Phân tích nguyên nhân vì sao khóa ECC kích thước 256-bit có thể cung cấp độ an toàn tương đương khóa RSA 3072-bit?
  \item Nếu kẻ tấn công sở hữu máy tính lượng tử có khả năng chạy thuật toán Shor, thuật toán mã hóa đối xứng AES-256 có bị bẻ gãy hoàn toàn như RSA hay không? Giải thích.
\end{enumerate}
\end{tukiemtra}

\begin{kbox}
Toán học mật mã tạo ra các cặp khóa nhưng bất lực trước tấn công mạo danh (MitM). Hạ tầng khóa công khai (PKI) là hệ thống kỹ thuật và thể chế bắt buộc để ràng buộc khóa công khai với định danh pháp lý (Identity Binding). Kỹ thuật RSA-CRT tăng tốc xử lý ký số lên 4 lần. ECC vượt trội RSA về hiệu năng và kích thước khóa ở mức an toàn 128-bit trở lên.
\end{kbox}
